The generation condition can also be restricted to an idempotent ideal.
The largest-idempotent-subideal construction is ordinary ideal theory;
for a unital ring, see Emmanouil \cite[Lemma~1.3(ii)]{Emmanouil}.
We give the small-category argument to specify its application here.

\begin{proposition}[the idempotent part of the canonical ideal]
\label{prop:canonical-ideal-core}
For a small preadditive category $\mathcal A$ and a two-sided ideal
$J$, there is a largest idempotent two-sided ideal $C(J)\subseteq J$.
It is the stable value of the transfinite sequence
\[
 J_0=J,\qquad J_{\alpha+1}=J_\alpha^2,\qquad
 J_\lambda=\bigcap_{\alpha<\lambda}J_\alpha
 \quad(\lambda\text{ a limit ordinal}).
\]
For every finitely generated projective covariant $\mathcal A$-module
$P$,
\[
 JP=P\quad\Longleftrightarrow\quad C(J)P=P.
\]
Every finite idempotent matrix with typed entries in $J$ already has
all its entries in $C(J)$. Thus Theorem~\ref{thm:canonical-ideal}
remains unchanged if $\mathcal J_L$ is replaced by
$C(\mathcal J_L)$.
\end{proposition}

\begin{proof}
Products of ideals consist of finite sums of typed composites;
arbitrary intersections and sums are ideals. Smallness ensures that
the family of ideals is a set. Let $C$ be the sum of all idempotent
ideals $I\subseteq J$. Each $I=I^2$ lies in $C^2$, so $C\subseteq
C^2$. The reverse inclusion follows because $C$ is an ideal.
Therefore $C$ is the largest idempotent ideal contained in $J$.

The displayed transfinite sequence is descending. Every idempotent
$I\subseteq J$ lies in each term: use $I=I^2$ at successors and
intersection at limits. The sequence stabilizes because the total
typed arrow collection of $\mathcal A$ is a set. Otherwise each
successor before its successor cardinal would remove a distinct
arrow, which is impossible. At a stable stage the ideal is idempotent,
so its maximality property identifies it with $C$. No finite-stage
stabilization has been assumed.

Let $E$ be a finite idempotent matrix with entries in $J$, and let
$I_E$ be the two-sided ideal generated by its entries. The equations
$E_{ij}=\sum_k E_{ik}E_{kj}$ put every generator in $I_E^2$.
Since $I_E^2$ is an ideal, $I_E\subseteq I_E^2$; hence
$I_E^2=I_E$. Maximality gives $I_E\subseteq C$, proving the
entrywise assertion.

For completeness, choose a finite representable sum $N$ and split
maps $P\xrightarrow{i}N\xrightarrow{r}P$ with $ri=1_P$.
If $JP=P$, then $i(P)=J i(P)\subseteq JN$, so all entries of the
idempotent $E=ir$ lie in $J$. They therefore lie in $C$, giving
$E(N)\subseteq CN$. Applying $r$ gives $P=rE(N)\subseteq CP$,
and consequently $CP=P$. The converse follows from $C\subseteq J$.
Apply this equivalence to the action category and its canonical ideal.
\end{proof}

This is square descent; it is not the iteration
$J_{\alpha+1}=J J_\alpha$ with a fixed left factor, whose stable
value need not be idempotent. The distinction is explicit in
\cite[Lemmas~1.1 and~1.3]{Emmanouil}. The already realized case
$J=2\mathbb Z$ has $C(J)=0$: its finite square-descending terms are
$2^{2^n}\mathbb Z$, with zero intersection at $\omega$. Thus the
reduction can be strict and need not stabilize at a finite stage.
If $C(\mathcal J_L)=0$, every allowed projective amplitude vanishes,
while its support $L$ is retained. No local-unit property, decision
procedure or classification of projectives over arbitrary idempotent
ideals is asserted.
